\section{Version 5: a falsified exploration hypothesis}
\subsection{Decision regret instead of trace count}
The inherited controller charges a worst-case loss but orders feasible probes by
how many observable traces they distinguish. This can reward distinctions that do
not change which program should be executed. We tested a replacement that values
uncertainty about the decision itself. For the fixed program library $\PP$, define
\begin{equation}
 \rho_g(\C)=\min_{p\in\PP}\max_{m\in\C}
 \left[\max_{q\in\PP}V_m(q,g)-V_m(p,g)\right].
\end{equation}
This is minimax simple regret restricted to the supplied library, not realized
cumulative regret or unrestricted optimal control. Model-wise optima are computed
inside the learner's compatible class; the true environment is not revealed.
Models with identical observable traces for every library program are quotiented
before scoring. Label multiplicity is therefore not treated as probability.

Let $\C_{u,\tau}$ be the nonempty cell consistent with observable trace $\tau$
after program $u$. The fixed two-objective score is
\begin{equation}
 G(u,\C)=\sum_g\rho_g(\C)-\max_{\tau}\sum_g\rho_g(\C_{u,\tau}).
 \label{eq:regretprobe}
\end{equation}
Among exactly checked, affordable probes, the new controller maximizes
$G/(1+d)$, then favors lower debit, larger optimistic current-objective gain,
and proposer order. The incumbent is eligible as a free probe. If the best
$G$ is zero, the controller executes the incumbent. Strict guaranteed
improvements still use the unchanged v4 admission rule. Enumeration has a strict
cap and returns \unk\ when incomplete. Exactness is conditional, exponential,
and limited to the fixed finite model and policy classes.

\subsection{Frozen protocol and rejected primary hypothesis}
The selection rule was recorded before a 20-seed pilot (20000--20019) and a
100-seed holdout (30000--30099). Every stream uses the inherited two-state,
two-action, two-output generator, 16 action words, horizon four, and 48-episode
A--B--A objective schedule. Four arms yield 3,840 pilot and 19,200 holdout rows:
v4 B16, decision-regret B16, its untrained ablation, and the original optimistic
full-history planner. All receive their own feedback and identical interaction
horizons. Total CPU time is measured but unequal. The primary contrast is
all-episode mean reward of the new trained controller minus v4 B16; intervals
pair independent seeds rather than treating episodes as independent.

\begin{table*}[t]
\centering\begin{tabular}{lrrrr}
\toprule
System & Mean return & Late A & Max. deficit & Seconds/stream \\
\midrule
Witness v4, B16 & 5.9217 & 5.8000 & 4 & 0.0359 \\
Regret probes, B16 & 5.8271 & 5.7200 & 10 & 0.0680 \\
Regret probes, no training & 5.8600 & 5.8000 & 10 & 0.0645 \\
Full-history planner & 5.9446 & 5.8000 & 4 & 0.0466 \\
\bottomrule
\end{tabular}

\caption{New frozen v5 holdout: 100 seeds, 48 episodes, four systems. Rewards
range from zero to eight. Deficit is the worst completed-prefix loss relative to
the initial anchor. Times cover measured per-stream decision/update mechanics,
not setup or a dedicated timing experiment; concurrent host load can affect them.}
\label{tab:v5}
\end{table*}

The proposed improvement is rejected. New mean return is 5.82708 versus
5.92167 for v4: difference $-0.09458$, paired 95\% $t$ interval
$[-0.16915,-0.02001]$. Against full-history planning the difference is
$-0.11750$, interval $[-0.19182,-0.04318]$. Late-A return falls from 5.80
to 5.72, with three seeds worse and 97 tied; its paired interval crosses zero.
The largest observed anchor deficit increases from 4 to 10, still within the
allowed 16. No protected-incumbent violations occur. Mean measured mechanics
cost is approximately 1.90 times v4. Fewer paired-check nodes do not imply lower
cost: the new arm adds about 5,957 model rollouts per stream for regret scoring.

Training also fails to show a clear benefit in this comparison: the trained
minus untrained difference is $-0.03292$, interval $[-0.09639,0.03055]$.
This secondary interval is descriptive and uncorrected for multiplicity.
The result does not contradict conditional incumbent retention; it shows that
retention alone neither chooses useful probes nor ensures efficient learning.
The pilot contains 20 distinct tables and the holdout 85, with 10 shared.
The new holdout also shares seven tables with v4 development and 22 with its
holdout. These are samples with replacement from only 256 labelled tables, not
disjoint environments or evidence of new-domain transfer.

\begin{figure*}[t]
\centering\includegraphics[width=.98\textwidth]{figures/v5_learning.pdf}
\caption{Frozen v5 result from raw episode records. Left: mean cumulative reward,
including all exploratory episodes. Right: per-seed mean differences sorted for
visibility; the seed is the statistical unit. The public objective changes at
16 and 32 completed episodes. Underlying CSVs accompany the figure.}
\label{fig:v5}
\end{figure*}

\subsection{Complementary free probes defeat one-step scoring}
An independent post-hoc counterexample isolates a completeness failure.
There are four compatible worlds indexed by hidden bits $(a,b)$. An anchor
returns 5 and reveals nothing. Two other actions also return 5 but reveal $a$
and $b$, respectively. Two terminal decision actions return 10 or 0 depending
on the parity $a\mathbin{\mathrm{xor}}b$. Observations encode these outcomes;
all four worlds are valid stationary reset models in a separately declared
finite class. This construction uses more outputs/actions than the primary
binary study and is identified as such in its executable artifact.

Initially the minimax regret is 5. Either free bit probe leaves parity ambiguous,
so its immediate regret reduction is zero. The decision actions have positive
regret reduction but a worst-case debit of 5 and are infeasible at budget zero.
The greedy selector therefore repeats the anchor indefinitely. Nevertheless,
executing the two free probes identifies parity and permits a guaranteed
improvement of 5 without any exploration loss. The reproducer executes both the
stall and the successful probe sequence; it is a finite implementation test and
written induction, not an additional Lean theorem.

\paragraph{A separately diagnosed admission plateau.}
Three frozen-holdout seeds stall after the sum of minimax regrets reaches zero.
A common optimal candidate exists, but it ties the incumbent in at least one
compatible model, so the strict $L>0$ promotion rule refuses it. Every probe also
has zero regret-reduction score. A separate weak-dominance closure promotes a
candidate only when exact comparison gives $L\geq0$ and $U>0$, preserving the
conditional retention premise. This is a post-hoc repair, not the frozen method.

An exhaustive development diagnostic evaluates all 256 labelled tables under
one proposer initialization per table, four arms and 48 episodes: 49,152 rows.
The repair improves the original v5 mean by 0.05453, improves ten tables and
worsens none, and repairs four late-return deficits in that diagnostic. It still
trails v4 and the full-history planner (mean 6.25326 versus 6.26123 and 6.29541)
and is slower. This finite schedule census supports the diagnosed repair only;
it supplies no held-out confidence interval or guarantee for other schedules,
model classes, or random initializations. It does not resolve the complementary
free-probe obstruction above.

\paragraph{Successor hypothesis.}
A bounded, adaptive two-probe plan can value complementary information that
Equation~\eqref{eq:regretprobe} misses. For each first outcome, choose a second
probe using that branch's updated cover, and maximize the worst-branch terminal
reduction. Charge every executed probe separately using the original ledger;
simulated outcomes supply planning information, never training evidence. A
cheaper baseline simply falls back to v4's trace-count rule when $G=0$.
Neither approach is entitled to a reward guarantee from uncertainty contraction.
The first kill test is failure to solve the complementary-probe example at zero
budget; the substantive kill test is no future reward advantage over both v4
and the cheaper fallback at equal total compute. The existing holdout is now
consumed for design purposes. Any tuned successor evaluation must be labelled
post-hoc or use a genuinely new task family/protocol.

\section{Actual local-model inference and native integration}
A first actual local-model smoke run used the already available 27B Q3 GGUF model
behind a loopback llama.cpp server, with temperature zero, a 96-token output cap,
and one call per episode for each model arm. The arms are raw executed-history
prompting, the same prompt plus the v4 checker and online proposer, and a
checker-based static-library control with no model calls. Programs are frozen
before execution; invalid output consumes the call and triggers an explicitly recorded fallback,
with the rejection retained. No offline retraining or paid API is used.

The first fixed seed, 41000, produced 36 episodes and 24 model calls with no
invalid outputs. All arms returned 5.3333 on average. Evaluator-only enumeration
then found every program had identical reward in that world, so this result has
no learning headroom. It establishes only executable inference and accounting.
A subsequent four-consecutive-seed exploratory extension (41001--41004)
produced 144 episodes and 96 model calls. Mean returns were 3.9583 for raw
history, 5.5417 for checked history, and 5.0000 for the static-library control.
Both model arms used 13,540 input and 576 output tokens, with respectively
141.91 and 128.51 seconds of observed model-call latency. Both stages together
contain 120 actual calls, all with valid output and complete usage records.
The static-library control still updates the small auxiliary ranker; only its
candidate library and absence of LLM calls are static.

These are five environments, not 60 independent trials per arm. No stateless
LLM control, native task, or powered retention evaluation ran. The symbolic
checker has a strong finite-class prior, more persistent state, extra CPU work,
and online auxiliary training. No LLM weights change. The pilot supports
integration feasibility and motivates a native experiment; it does not isolate
memory quality, parameter learning, or cost-matched superiority. Initial model
loading took about 77.7 seconds and is excluded from call timings. Host/GPU
memory values are environment observations, not isolated peak-memory benchmarks.

An independent review found misleading malformed-output metadata and incomplete
usage-field handling in the pilot driver. Neither branch was exercised in these
runs: every completion was valid and every usage record complete. Future-run
code and regression tests are corrected; the exact executed source, original
results, and review are preserved. For invalid future output, raw history uses
the anchor while the checked arm skips registration and uses its existing
controller, with no extra call. This distinction is now explicit.


The native integration targets the official CL-Bench interface at commit
\texttt{5f8c50eb1e84b2eda2ef4faff757dfc812a0ea26}. Its feedback semantics are
material: the runner separately delivers observations; duplicating them through
query handling would change the ICL baseline. A task may legally reveal a correct
answer in post-action feedback; that is different from reading the evaluator's
answer records before acting. Native reward and each system's stateful--stateless
gain must remain separate. A bridge and interface tests do not supply a native
benchmark result, an empirical retention result, or a valid finite-state
abstraction for database drift. Those remain the next experimental obligations.


Twenty optional-runtime tests pass, including eight tests executing the actual
pinned upstream ICL/runner and SQLite task; twelve dependency-free contracts also
belong to the 401-test standard suite. The native transport is scripted for these
checks, and the SQLite database/questions are fixtures. Its exact certificate
metadata is always \unk. The bridge is a baseline and audit foundation, not an
implemented native Witness adaptation algorithm. The unmodified official
benchmark dataset and a live native model transport remain unexecuted.
